\subsection{Demand}\label{sec:demand}
\tbd

\subsubsection{Adjustments}\label{sec:demand-adjust}

\paragraph{Place demand inside a country with published shares} (\commit{563658d5})

\begin{figure}[htb]
  \centering
  \includegraphics[width=\textwidth]{demand_placement.pdf}
  \caption{Electricity demand per province of China in the model over its published share of national consumption in 2024, with the province's published share next to its name.
  Grey: the load distributed by GDP and population.
  Blue: with the provincial shares.
  Demand is counted where the buses of the base network lie.}
  \label{fig:demandplacement}
\end{figure}

\textit{What needed fixing.}
PyPSA-Earth distributes a country's hourly load to the buses with one weight per first-level administrative unit, 0.6 times its share of GDP plus 0.4 times its share of population, a regression on European countries.
Inside a unit the load follows the area of the bus regions.
Where heavy industry sits in units with few people, the load ends up in the wrong place.
In China, 15\,\% of national demand lay in the wrong province: Inner Mongolia, Xinjiang, Ningxia and Qinghai had a third of their published consumption, Heilongjiang and Jilin twice theirs (Figure~\ref{fig:demandplacement}).
The model then moved the missing demand's generation east: Inner Mongolia exported \SI{711}{TWh} on a load of \SI{188}{TWh}.

\textit{What was done.}
A table gives the share of the country's demand per administrative unit.
A unit with a row gets that share of the country's load in every hour; units without a row keep the weight from GDP and population, and the table's shares are scaled to what is left.
Inside a unit the load is distributed as before, and a country without rows is not touched.
The national level stays the calibrated one, and every planning horizon uses the same shares.
For China the table holds the consumption of the 31 provinces in 2024, Hong Kong and Macau.
No table of 2024 could be read as text, so the provinces are measured off a published bar chart of the China Electricity Council's figures and scaled to the national total; for the eight provinces with an exact figure from another source the reading is within 2\,\%.
The feature becomes active in rule \texttt{build\_demand\_profiles}, for China only:
\yamlblock{demand_shares}

\textit{Check.}
With the shares, 0.2\,\% of China's demand lies in another province than published; what is left comes from bus regions that reach across a provincial border.
In the dispatch of 2024 Inner Mongolia then exports \SI{358}{TWh} and Xinjiang \SI{108}{TWh}; their published generation minus consumption of 2024 is about 300 and \SI{105}{TWh}.
All provinces share one national hourly demand shape, and the shares do not change over the horizons.

\paragraph{Move the load of a part of a country in time} (\commit{8f290c58})

\textit{What needed fixing.}
The model has one hourly load series per country.
The two US regions are parts of one country and both got the national series, which blends four time zones.
The load of the West runs later in the day than this blend and that of the East a little earlier, so the West's evening peak fell too early against its solar output.

\textit{What was done.}
A table gives the hours by which the series of a country is moved in time before it is scaled and distributed.
A positive value moves it later; fractions of an hour are interpolated, the series wraps around at the ends of the year, and its sum is kept.
For each US region the shift is its mean offset to UTC against that of the contiguous United States, with the states weighted by their retail electricity sales of 2024: US West runs 1.77 hours later and US East 0.39 hours earlier.
A state that a time-zone border cuts is counted with its larger part.
The shape itself stays the national one: differences in the seasonal pattern and in the form of the daily curve between the regions are not represented.
The feature becomes active in rule \texttt{build\_demand\_profiles}, for the two US regions:
\yamlblock{time_shift}

\textit{Check.}
In US West the mean daily peak of the load moves from 23:24 to 01:16 UTC, from 15:24 to 17:16 Pacific Standard Time, and the annual demand is unchanged.
